\documentclass[../../../include/open-logic-section]{subfiles}

\begin{document}

\olfileid[fr]{sth}{z}{story}
\olsection{La formation des ensembles, plus en détail}
	
Dans la \olref[story][approach]{sec}, nous avons cité la description
que donne Shoenfield de la formation des ensembles. Nous voulons
maintenant énoncer quelques principes supplémentaires pour préciser
ce récit. Les voici :
\begin{enumerate}
\item[] \stageshier. Tout ensemble se forme à une étape.
\item[] \stagesord. Les étapes sont ordonnées : certaines en
\emph{précèdent} d'autres.\footnote{Nous supposerons en réalité,
tacitement, que les étapes sont \emph{bien ordonnées}. La
\olref[ordinals][]{chap} expliquera ce que cela signifie. Il s'agit
d'une hypothèse substantielle. En fait, grâce à une technique très
ingénieuse due à \citet{Scott1974}, on peut \emph{se passer} de cette
hypothèse et ensuite la \emph{démontrer}. (Cela expliquera aussi
pourquoi nous pouvons penser qu'il existe une étape initiale.) Nous
ne pouvons développer ce point ici ; pour en savoir plus, voir
\citet{ButtonLT1}.}
\item[] \stagesacc. Pour toute étape $S$ et toute collection
d'ensembles formés \emph{avant} l'étape~$S$, il se forme à l'étape~$S$
un ensemble dont les membres sont exactement ceux de la collection.
Aucun autre ensemble ne se forme à l'étape~$S$.
\end{enumerate}
Ces principes sont informels, mais ils nous permettront de justifier
plusieurs axiomes de la théorie des ensembles de Zermelo.

(Une mise en garde s'impose. Nous présenterons certains axiomes tout
à fait classiques, sous leurs noms habituels, mais les principes en
italique que nous venons d'énoncer n'ont pas de noms établis dans la
littérature. Nous leur donnons simplement des noms qui, nous
l'espérons, seront utiles.)

\end{document}
